\documentclass[11pt]{article}
\usepackage{mathunal-base}
\mathunalfuente{serif}
\usepackage{amssymb}

\newcommand{\resp}{\underline{\hspace{2.5cm}}}
\newcommand{\resplong}{\underline{\hspace{9cm}}}

\newenvironment{pseudo}{%
  \begin{quote}\ttfamily\small\begin{tabbing}
  \hspace{2em}\=\hspace{2em}\=\hspace{2em}\=\hspace{2em}\=\kill}{%
  \end{tabbing}\end{quote}}

\begin{document}

{\large\bfseries Primer Parcial de Matemáticas Discretas --- 16 de septiembre de 2023}

\medskip
\noindent Escuela de Matemáticas. Universidad Nacional de Colombia, Sede Medellín.

\medskip
\noindent\textit{Este documento reúne los dos temas del examen (A y B). Duración: 1 hora y 50 minutos. Seis preguntas. En las preguntas 1 a 4 conteste en el espacio destinado para cada respuesta; la calificación tendrá en cuenta sólo la respuesta. En las preguntas 5 y 6 muestre detalladamente el procedimiento.}

\medskip
\noindent Puntaje: 1 (16\,\%), 2 (12\,\%), 3 (20\,\%), 4 (12\,\%), 5 (20\,\%), 6 (20\,\%).

\bigskip\hrule\medskip
\begin{center}\bfseries\large Tema A\end{center}

\begin{enumerate}
\item (16\,\%) Sean $p$, $q$ y $r$ proposiciones. Complete los valores indicados en la siguiente tabla de verdad (las casillas con porcentaje son las que debe completar):
\begin{center}
\begin{tabular}{|c|c|c|c|c|c|c|}
\hline
$p$ & $q$ & $r$ & $p \to q$ & $\neg r$ & $(\neg r) \wedge (p \to q)$ & $p \leftrightarrow q$ \\
\small(4\,\%) & \small(4\,\%) & & \small(2\,\%) & \small(3\,\%) & \small(3\,\%) & \\
\hline
 & & V & & & & F \\
\hline
\end{tabular}
\end{center}

\item (12\,\%) Considere las proposiciones $p$, $q$ dadas por
\[
p : \ \forall x\ \exists y\ (x \leq y^2 + 1), \qquad
q : \ \exists z\ \forall w\ (z + 1 = w),
\]
donde el dominio de discurso son los números enteros ($x,y,z,w \in \mathbb{Z}$). Determine la contrarrecíproca (o contrapositiva) de la afirmación $p \to q$. (Puede emplear los símbolos $\neq$, $<$, $\leq$, $>$, $\geq$.)

\vspace{4pt}
Respuesta: \resplong

\item (20\,\%)
\begin{enumerate}
\item Un pequeño robot está programado para dar pasos en línea recta de longitudes $2$ ó $3$ centímetros. Si $P(n)$ es el número de formas de recorrer un camino de $n$ centímetros, note que $P(1) = 0$ y $P(2) = P(3) = 1$. Complete el siguiente pseudocódigo recursivo para calcular la función $P(n)$:
\begin{pseudo}
Procedure P($n$)\\
\>If $n == 1$, Return $0$\\
\>If $n == 2$ or $n == 3$, Return \underline{\hspace{2cm}} \quad\textbf{(2\,\%)}\\
\>Else\\
\>\>Return \underline{\hspace{4cm}} \quad\textbf{(10\,\%)}\\
End
\end{pseudo}
\emph{Nota:} las palabras \texttt{Procedure} y \texttt{End} indican dónde comienza y termina el pseudocódigo.
\item (8\,\%) Se da el siguiente código para calcular recursivamente la sucesión $P_a(n)$:
\begin{pseudo}
Procedure Pa($n$)\\
\>If $n == 1$, Return $8$\\
\>If $n == 2$, Return $4$\\
\>Else\\
\>\>Return $6 \cdot P_a(n-1) - 2 \cdot P_a(n-2)$\\
End
\end{pseudo}
Si damos el valor $n = 5$, el valor que retorna este programa es $P_a(5) = \resp$.
\end{enumerate}

\item (12\,\%) Se requiere un código que calcule la suma de las entradas mayores a $100$ de una lista $s = [s_1, s_2, \dots, s_n]$ de números enteros. Si la lista no tiene valores mayores que $100$, la salida debe ser cero. Por ejemplo, si $s = [100, 65, 101, -1, 110]$, el programa debe retornar $101 + 110 = 211$. Complete el siguiente código recursivo:
\begin{pseudo}
\textbf{Entrada:} la lista $s = [s_1, \dots, s_n]$ y la longitud $n$.\\
\textbf{Salida:} la suma de los elementos mayores a $100$, ó $0$ si no hay tales elementos.\\[2pt]
Mayores\_a\_100($s$, $n$)\\
\>If $n == 0$, Return \underline{\hspace{2cm}} \quad\textbf{(2\,\%)}\\
\>$x = $ Mayores\_a\_100($s$, $n-1$)\\
\>If \underline{\hspace{3cm}} \quad\textbf{(4\,\%)}\\
\>\>Return \underline{\hspace{3cm}} \quad\textbf{(4\,\%)}\\
\>Else, Return \underline{\hspace{2cm}} \quad\textbf{(2\,\%)}
\end{pseudo}
\emph{Nota:} si $n = 0$ la lista es vacía. Además, Mayores\_a\_100($s$, $n-1$) indica la función aplicada a la lista $[s_1, \dots, s_{n-1}]$ de longitud $n-1$.

\item (20\,\%) Demuestre que si $n$ es un número entero, entonces $20 \cdot n + 2023$ es impar.

\item (20\,\%) Demuestre empleando inducción matemática que la siguiente afirmación es válida para todo entero $n \geq 1$:
\[
(3 \cdot 5) + (3 \cdot 5^2) + \cdots + (3 \cdot 5^n) = \frac{15}{4}(5^n - 1).
\]
\end{enumerate}

\bigskip\hrule\medskip
\begin{center}\bfseries\large Tema B\end{center}

\begin{enumerate}
\item (16\,\%) Sean $p$, $q$ y $r$ proposiciones. Complete los valores indicados en la siguiente tabla de verdad (las casillas con porcentaje son las que debe completar):
\begin{center}
\begin{tabular}{|c|c|c|c|c|c|c|}
\hline
$p$ & $q$ & $r$ & $q \to p$ & $\neg r$ & $(\neg r) \vee (q \to p)$ & $q \leftrightarrow p$ \\
\small(4\,\%) & \small(4\,\%) & & \small(2\,\%) & \small(3\,\%) & \small(3\,\%) & \\
\hline
 & & F & & & & F \\
\hline
\end{tabular}
\end{center}

\item (12\,\%) Considere las proposiciones $p$, $q$ dadas por
\[
p : \ \forall x\ \exists y\ (x^2 \geq y + 1), \qquad
q : \ \exists z\ \forall w\ (z + 3 = w),
\]
donde el dominio de discurso son los números enteros ($x,y,z,w \in \mathbb{Z}$). Determine la contrarrecíproca (o contrapositiva) de la afirmación $p \to q$. (Puede emplear los símbolos $\neq$, $<$, $\leq$, $>$, $\geq$.)

\vspace{4pt}
Respuesta: \resplong

\item (20\,\%)
\begin{enumerate}
\item Un pequeño robot está programado para dar pasos en línea recta de longitudes $1$ ó $3$ centímetros. Si $Q(n)$ es el número de formas de recorrer un camino de $n \geq 1$ centímetros, note que $Q(1) = Q(2) = 1$ y $Q(3) = 2$. Complete el siguiente pseudocódigo recursivo para calcular la función $Q(n)$:
\begin{pseudo}
Procedure Q($n$)\\
\>If $n == 1$ or $n == 2$, Return $1$\\
\>If $n == 3$, Return \underline{\hspace{2cm}} \quad\textbf{(2\,\%)}\\
\>Else\\
\>\>Return \underline{\hspace{4cm}} \quad\textbf{(10\,\%)}\\
End
\end{pseudo}
\emph{Nota:} las palabras \texttt{Procedure} y \texttt{End} indican dónde comienza y termina el pseudocódigo.
\item (8\,\%) Se da el siguiente código para calcular recursivamente la sucesión $T_a(n)$:
\begin{pseudo}
Procedure Ta($n$)\\
\>If $n == 1$, Return $6$\\
\>If $n == 2$, Return $4$\\
\>Else\\
\>\>Return $8 \cdot T_a(n-1) - 4 \cdot T_a(n-2)$\\
End
\end{pseudo}
Si damos el valor $n = 5$, el valor que retorna este programa es $T_a(5) = \resp$.
\end{enumerate}

\item (12\,\%) Se requiere un código que calcule la suma de las entradas menores a $50$ de una lista $s = [s_1, s_2, \dots, s_n]$ de números enteros. Si la lista no tiene valores menores que $50$, la salida debe ser cero. Por ejemplo, si $s = [20, 65, 32, 50, -1]$, el programa debe retornar $20 + 32 + (-1) = 51$. Complete el siguiente código recursivo:
\begin{pseudo}
\textbf{Entrada:} la lista $s = [s_1, \dots, s_n]$ y la longitud $n$.\\
\textbf{Salida:} la suma de los elementos menores a $50$, ó $0$ si no hay tales elementos.\\[2pt]
Menores\_a\_50($s$, $n$)\\
\>If $n == 0$, Return \underline{\hspace{2cm}} \quad\textbf{(2\,\%)}\\
\>$x = $ Menores\_a\_50($s$, $n-1$)\\
\>If \underline{\hspace{3cm}} \quad\textbf{(4\,\%)}\\
\>\>Return \underline{\hspace{3cm}} \quad\textbf{(4\,\%)}\\
\>Else, Return \underline{\hspace{2cm}} \quad\textbf{(2\,\%)}
\end{pseudo}
\emph{Nota:} si $n = 0$ la lista es vacía. Además, Menores\_a\_50($s$, $n-1$) indica la función aplicada a la lista $[s_1, \dots, s_{n-1}]$ de longitud $n-1$.

\item (20\,\%) Demuestre que si $n$ es un número entero, entonces $14 \cdot n + 2025$ es impar.

\item (20\,\%) Demuestre empleando inducción matemática que la siguiente afirmación es válida para todo entero $n \geq 1$:
\[
(5 \cdot 3) + (5 \cdot 3^2) + \cdots + (5 \cdot 3^n) = \frac{15}{2}(3^n - 1).
\]
\end{enumerate}

\vfill
\noindent{\small\itshape Transcripción digital --- MathUNAL}

\end{document}
